\chapter{Factorized Scattering and Integrability}
\label{ch:integrable-factorized-scattering}

\ConventionsLorentzian

This volume begins with massive two-dimensional QFT because that is the place
where integrability has a sharp scattering-theoretic meaning.  Let the theory
have a vacuum representation on \(\mathbb M^2\), a mass gap above the vacuum,
and isolated one-particle species labelled by \(a,b,\ldots\), with masses
\(m_a>0\).  A one-particle momentum is parametrized by rapidity
\[
  p_a^\mu(\theta)
  =
  (m_a\cosh\theta,\ m_a\sinh\theta),
  \qquad \theta\in\R .
\]
The Lorentz group acts by translating \(\theta\).  Hence a two-particle
invariant depends on the rapidity difference \(\theta_{12}=\theta_1-\theta_2\).

\section{Factorized Scattering Datum}

\begin{definition}[Elastic factorized scattering datum]
\label{def:elastic-factorized-scattering-datum}
An elastic factorized scattering datum consists of a finite or countable set
of massive species, a positive mass \(m_a\) for each species, and meromorphic
functions
\[
  S_{ab}^{cd}(\theta),
  \qquad a,b,c,d\in I,
\]
defined on a common rapidity-domain cover, satisfying the following algebraic
and analytic conditions.

First, two-particle unitarity is
\[
  \sum_{e,f}
  S_{ab}^{ef}(\theta)\,
  S_{ef}^{cd}(-\theta)
  =
  \delta_a^{c}\delta_b^{d}.
\]
Second, Hermitian analyticity is
\[
  S_{ab}^{cd}(\theta)^\ast
  =
  S_{cd}^{ab}(-\theta^\ast)
\]
on the boundary where both sides are defined.  Third, if
\(\bar a\) denotes the charge-conjugate species, crossing has the form
\[
  S_{ab}^{cd}(\theta)
  =
  C_{a\bar a'}C^{c\bar c'}\,
  S_{\bar c' b}^{\bar a' d}(\ii\pi-\theta),
\]
with the charge-conjugation matrix \(C\) fixed as part of the particle-state
normalization.  Fourth, three-particle consistency is the Yang--Baxter
equation
\[
  S_{12}(\theta_{12})S_{13}(\theta_{13})S_{23}(\theta_{23})
  =
  S_{23}(\theta_{23})S_{13}(\theta_{13})S_{12}(\theta_{12})
\]
as an identity on the tensor product of three one-particle internal spaces.
\end{definition}

The notation \(S_{ij}\) means that the two-particle matrix \(S\) acts on the
\(i\)-th and \(j\)-th tensor factors and as the identity on the remaining
factor.  The Yang--Baxter equation is the associativity condition for
reordering three rapidities by adjacent transpositions.  It is the exact
analogue, for ordered asymptotic particles in \(1+1\) dimensions, of the
ordinary associativity constraint on multiple OPE expansions in the CFT
volume.

\begin{definition}[Separating higher-spin family]
\label{def:separating-higher-spin-family}
Let \(\{Q_{s,\lambda}\}_{(s,\lambda)\in\mathcal A}\) be commuting conserved
charges whose one-particle eigenvalues have the form
\[
  q_{s,\lambda,a}(\theta)=\kappa_{s,\lambda,a}\,\ee^{s\theta},
  \qquad s>1.
\]
For fixed particle number \(N\), set \(z_i=\ee^{\theta_i}\).  The family is
\(N\)-separating on a rapidity domain \(U_N\) if the moment map
\[
  \Phi_N\bigl(\{(a_i,z_i)\}_{i=1}^N\bigr)
  =
  \left(
    \sum_{i=1}^N\kappa_{s,\lambda,a_i}z_i^s
  \right)_{(s,\lambda)\in\mathcal A}
\]
is injective on \(U_N/\mathfrak S_N\), outside a specified closed exceptional
set consisting of coincident rapidities and declared threshold or bound-state
singular loci.  The quotient means that only the unordered multiset of
species and rapidities is to be recovered.

A concrete sufficient condition, for a finite species set, is the following.
For each species \(a\) and for each \(r\) needed up to the maximal particle
number under consideration, finite linear combinations of the charges produce
the species-resolved power sums
\[
  p_{r,a}=\sum_{i:\,a_i=a}z_i^r .
\]
Then Newton identities recover the monic polynomial
\[
  \prod_{i:\,a_i=a}(t-z_i)
\]
for each species \(a\), and hence recover the unordered species-resolved
rapidity multiset wherever the roots are distinct.  The definition above is
the hypothesis used below; this sufficient condition is only a convenient
checkable way to obtain it.
\end{definition}

\begin{proposition}[Factorization from conserved higher-spin charges]
\label{prop:higher-spin-factorization-kinematic}
Assume a massive \(1+1\)-dimensional QFT has Haag--Ruelle asymptotic particle
states, a family of conserved charges
\(\{Q_{s,\lambda}\}_{(s,\lambda)\in\mathcal A}\) as in
Definition~\ref{def:separating-higher-spin-family}, and scattering kernels
which, after smearing by compactly supported wave packets in rapidity
chambers, depend continuously on the wave packets away from threshold and
bound-state singular loci.  Equivalently, the unsmeared kernels are
distributional boundary values whose singular support in the rapidity
variables is contained in those declared loci.

If the scattering operator commutes with every \(Q_{s,\lambda}\), and if the
charge family is \(N\)-separating on the incoming and outgoing \(N\)-particle
domains under consideration, then the \(N\to N\) scattering kernel is
supported, away from the declared singular loci, on the union of permutation
graphs
\[
  \{(\theta'_1,b_1;\ldots;\theta'_N,b_N)
    =
    (\theta_{\sigma(1)},a_{\sigma(1)};\ldots;
     \theta_{\sigma(N)},a_{\sigma(N)})\}.
\]
If the corresponding moment maps for \(M\neq N\) have disjoint generic
images, then generic \(N\to M\) particle production is absent.
\end{proposition}

\begin{proof}
Let \(K_{N\to M}\) be the distribution kernel of the scattering operator
between ordered rapidity chambers, away from the declared singular loci.
Commutation with \(Q_{s,\lambda}\) gives, in the sense of distributions,
\[
  \left(
    \sum_{i=1}^N\kappa_{s,\lambda,a_i}\ee^{s\theta_i}
    -
    \sum_{j=1}^M\kappa_{s,\lambda,b_j}\ee^{s\theta'_j}
  \right)
  K_{N\to M}=0
\]
for every \((s,\lambda)\).  A standard support property of distributions says
that if a smooth function \(f\) is nonzero on an open set and \(fT=0\), then
the distribution \(T\) vanishes on that open set.  Therefore the support of
the kernel is contained in the common zero locus of all differences of charge
eigenvalues.

For \(M=N\), the \(N\)-separation hypothesis says that this common zero locus,
outside the exceptional set, is precisely the union of permutation graphs
displayed in the statement.  Thus, after smearing with wave packets supported
in a connected component of the regular rapidity chamber, the scattering
operator cannot have support away from elastic permutations.  For \(M\neq N\),
the assumed disjointness of the generic moment-map images leaves no regular
common zero locus, so the production kernel vanishes on the corresponding
open domain.  Threshold and bound-state loci are excluded from the statement
because the distribution kernels may have additional singular components
there; those components must be described after the relevant channel is
resolved.
\end{proof}

\section{Zamolodchikov--Faddeev Operators}

The factorized datum becomes an operator algebra on asymptotic wavefunctions.
For each species \(a\) introduce distributional creation operators
\(Z_a^\dagger(\theta)\) obeying
\begin{equation}
  Z_a^\dagger(\theta_1)Z_b^\dagger(\theta_2)
  =
  S_{ab}^{cd}(\theta_1-\theta_2)\,
  Z_d^\dagger(\theta_2)Z_c^\dagger(\theta_1),
  \qquad \theta_1>\theta_2 .
  \label{eq:zf-creation-exchange}
\end{equation}
Repeated species labels are summed.  Associativity of the product of three
creation operators is equivalent to the Yang--Baxter equation in
Definition~\ref{def:elastic-factorized-scattering-datum}; applying
\eqref{eq:zf-creation-exchange} to reorder
\(\theta_1>\theta_2>\theta_3\) along the two reduced words
\((12)(13)(23)\) and \((23)(13)(12)\) gives precisely the two sides of that
equation.

\section{Local Operators and Form Factors}

Let \(\mathcal O(x)\) be a local operator whose matrix elements between the
vacuum and ordered out-states exist as distributions.  Its \(n\)-particle form
factor is
\[
  F^{\mathcal O}_{a_1\ldots a_n}(\theta_1,\ldots,\theta_n)
  =
  \bra{\vac}\mathcal O(0)
  Z_{a_1}^\dagger(\theta_1)\cdots Z_{a_n}^\dagger(\theta_n)\ket{\vac},
  \qquad \theta_1>\cdots>\theta_n .
\]
Changing the ordered rapidity chamber changes the scattering basis by the
two-body \(S\)-matrix.  Thus the form-factor boundary values obey Watson
exchange equations,
\[
  F^{\mathcal O}_{\ldots a_i a_{i+1}\ldots}(\ldots,\theta_i,\theta_{i+1},\ldots)
  =
  S_{a_i a_{i+1}}^{bc}(\theta_i-\theta_{i+1})\,
  F^{\mathcal O}_{\ldots c b\ldots}(\ldots,\theta_{i+1},\theta_i,\ldots),
\]
together with crossing, cyclicity, and residue equations at bound-state
poles.  The exchange equation is a statement about boundary values in
neighboring rapidity chambers and the normalization of ordered scattering
states; locality of \(\mathcal O\) enters the cyclicity and semi-locality
conditions developed in Chapter~\ref{ch:integrable-form-factor-bootstrap-local-operators}.
The operator reconstruction problem is to determine which solutions of these
functional equations come from genuine local operators in a Hilbert-space
QFT.  The theorem-boundary data include the local operator domain,
positivity, convergence of form-factor expansions, and clustering.

\begin{openproblem}[Completeness of form-factor reconstruction]
\label{op:form-factor-reconstruction-completeness}
Give intrinsic hypotheses on a factorized scattering datum and a family of
form-factor solutions that imply the existence of a local QFT with those
asymptotic particles, those local operators, and Wightman functions obtained
by convergent form-factor sums.  Existing constructive approaches prove this
for important classes of two-dimensional models; the monograph will develop
the exact hypotheses and examples in later chapters of this volume.
\end{openproblem}
